% ============================================================================
%  Rapport de stage PFA — Application mobile EMC Helpline
%  Sayouba Ouédraogo · CMRPI · 15 juillet – 31 août 2026
%
%  Compilation :  latexmk -xelatex rapport.tex
%  (ou :  xelatex rapport → biber rapport → xelatex ×2)
% ============================================================================
\documentclass[11pt,a4paper,oneside,openany]{report}
% ============================================================================
%  preambule.tex — configuration typographique, couleurs, boîtes, diagrammes
%  Compilation : XeLaTeX (ou LuaLaTeX) + Biber
% ============================================================================

% ---------------------------------------------------------------- géométrie
\usepackage[a4paper,top=2.3cm,bottom=2.2cm,left=2.4cm,right=2.4cm,
            headheight=14pt,footskip=1.0cm]{geometry}

% ---------------------------------------------------------------- polices
\usepackage{fontspec}
\usepackage{unicode-math}
\defaultfontfeatures{Ligatures=TeX}

% Chargement défensif : la police de secours est celle que toute installation
% TeX possède, ce qui évite qu'un rapport cesse de compiler ailleurs.
\IfFontExistsTF{TeX Gyre Pagella}
  {\setmainfont{TeX Gyre Pagella}}
  {\setmainfont{Latin Modern Roman}}
\IfFontExistsTF{TeX Gyre Heros}
  {\setsansfont{TeX Gyre Heros}[Scale=MatchLowercase]}
  {\setsansfont{Latin Modern Sans}[Scale=MatchLowercase]}
\IfFontExistsTF{DejaVu Sans Mono}
  {\setmonofont{DejaVu Sans Mono}[Scale=0.84]}
  {\setmonofont{Latin Modern Mono}[Scale=MatchLowercase]}
\IfFontExistsTF{TeX Gyre Pagella Math}
  {\setmathfont{TeX Gyre Pagella Math}}
  {\setmathfont{Latin Modern Math}}

\usepackage[french]{babel}
\frenchsetup{SmallCapsFigTabCaptions=false}
\usepackage{csquotes}
\usepackage{microtype}

% ---------------------------------------------------------------- couleurs
% Reprises de lib/core/constants/app_colors.dart : le rapport porte l'identité
% visuelle du produit qu'il décrit.
\usepackage{xcolor}
\definecolor{emcblue}{HTML}{1E3A8A}   % primaryBlue
\definecolor{emcorange}{HTML}{C2410C} % primaryOrange (AA sur blanc)
\definecolor{emcred}{HTML}{B91C1C}    % dangerRedStrong
\definecolor{emcslate}{HTML}{0F172A}  % textPrimary
\definecolor{emcgrey}{HTML}{64748B}   % textSecondary
\definecolor{emcborder}{HTML}{E2E8F0}
\definecolor{emcbg}{HTML}{F8FAFC}
\definecolor{emcgreen}{HTML}{15803D}

% ---------------------------------------------------------------- tableaux
\usepackage{booktabs}
\usepackage{tabularx}
\usepackage{multirow}
\usepackage{array}
\usepackage{longtable}
\newcolumntype{L}[1]{>{\raggedright\arraybackslash}p{#1}}
\newcolumntype{C}[1]{>{\centering\arraybackslash}p{#1}}
\renewcommand{\arraystretch}{1.18}
\usepackage[table]{colortbl}

% ---------------------------------------------------------------- listes
\usepackage{enumitem}
\setlist[itemize]{leftmargin=1.15em,itemsep=1pt,topsep=2pt,parsep=0pt}
\setlist[enumerate]{leftmargin=1.5em,itemsep=1pt,topsep=2pt,parsep=0pt}

% ---------------------------------------------------------------- figures
\usepackage{graphicx}
\usepackage{float}
% Les flottants remplissent les pages ; sans cela, un [H] trop haut laisse
% une demi-page blanche derrière lui.
\renewcommand{\topfraction}{0.92}
\renewcommand{\bottomfraction}{0.7}
\renewcommand{\textfraction}{0.06}
\renewcommand{\floatpagefraction}{0.75}
\setcounter{topnumber}{3}
\setcounter{bottomnumber}{2}
\setcounter{totalnumber}{5}
% Espacement des flottants resserré : le rapport tient en quinze pages.
\setlength{\intextsep}{7pt plus 2pt minus 2pt}
\setlength{\textfloatsep}{8pt plus 2pt minus 2pt}
\setlength{\floatsep}{7pt plus 2pt minus 2pt}
\setlength{\abovecaptionskip}{4pt}
\setlength{\belowcaptionskip}{0pt}
\usepackage[font=small,labelfont={bf,color=emcblue},
            justification=centering,skip=6pt]{caption}
\usepackage{subcaption}
\captionsetup[sub]{font=footnotesize,labelfont=bf}

% ---------------------------------------------------------------- TikZ
\usepackage{tikz}
\usetikzlibrary{shapes.geometric,shapes.multipart,shapes.symbols,arrows.meta,
                positioning,calc,fit,backgrounds,matrix,chains,
                decorations.pathreplacing,decorations.markings,patterns,quotes}
\usepackage{pgfgantt}

% ---------------------------------------------------------------- boîtes
\usepackage[most]{tcolorbox}
\tcbuselibrary{listings,breakable,skins}

\newtcolorbox{decision}[1][]{%
  enhanced,breakable,colback=emcblue!4,colframe=emcblue,
  boxrule=0pt,leftrule=2.5pt,arc=2pt,left=6pt,right=6pt,top=5pt,bottom=5pt,
  fonttitle=\sffamily\bfseries\small,coltitle=emcblue,
  attach boxed title to top left={yshift=-2mm,xshift=6mm},
  boxed title style={colback=white,boxrule=0pt,arc=1pt},
  title={Décision technique},#1}

\newtcolorbox{vigilance}[1][]{%
  enhanced,breakable,colback=emcorange!5,colframe=emcorange,
  boxrule=0pt,leftrule=2.5pt,arc=2pt,left=6pt,right=6pt,top=5pt,bottom=5pt,
  fonttitle=\sffamily\bfseries\small,coltitle=emcorange,
  attach boxed title to top left={yshift=-2mm,xshift=6mm},
  boxed title style={colback=white,boxrule=0pt,arc=1pt},
  title={Point de vigilance},#1}

\newtcolorbox{constat}[1][]{%
  enhanced,breakable,colback=emcbg,colframe=emcborder,
  boxrule=0.8pt,arc=3pt,left=7pt,right=7pt,top=5pt,bottom=5pt,#1}

% ---------------------------------------------------------------- code
\usepackage{listings}
\lstdefinelanguage{Dart}{
  morekeywords={abstract,as,assert,async,await,break,case,catch,class,const,
    continue,default,deferred,do,dynamic,else,enum,export,extends,external,
    factory,false,final,finally,for,get,if,implements,import,in,is,late,
    library,mixin,new,null,on,operator,part,required,rethrow,return,sealed,
    set,static,super,switch,sync,this,throw,true,try,typedef,var,void,while,
    with,yield,String,int,double,bool,List,Map,Future,Widget,record},
  sensitive=true,
  morecomment=[l]{//},morecomment=[s]{/*}{*/},morecomment=[l]{///},
  morestring=[b]',morestring=[b]"}
\lstdefinelanguage{TypeScript}{
  morekeywords={abstract,any,as,async,await,boolean,break,case,catch,class,
    const,constructor,continue,declare,default,delete,do,else,enum,export,
    extends,false,finally,for,from,function,get,if,implements,import,in,
    instanceof,interface,let,new,null,number,of,private,protected,public,
    readonly,return,set,static,string,super,switch,this,throw,true,try,type,
    typeof,undefined,var,void,while,yield},
  sensitive=true,
  morecomment=[l]{//},morecomment=[s]{/*}{*/},
  morestring=[b]',morestring=[b]",morestring=[b]`}

\lstdefinestyle{emc}{
  basicstyle=\ttfamily\scriptsize\color{emcslate},
  keywordstyle=\color{emcblue}\bfseries,
  commentstyle=\color{emcgrey}\itshape,
  stringstyle=\color{emcgreen},
  numbers=left,numberstyle=\ttfamily\tiny\color{emcgrey},numbersep=8pt,
  showstringspaces=false,breaklines=true,breakatwhitespace=true,
  tabsize=2,columns=fullflexible,keepspaces=true,
  frame=leftline,framerule=1pt,rulecolor=\color{emcborder},
  xleftmargin=14pt,framexleftmargin=14pt,
  aboveskip=8pt,belowskip=6pt,
  literate={é}{{\'e}}1 {è}{{\`e}}1 {ê}{{\^e}}1 {à}{{\`a}}1 {ç}{{\c c}}1
           {É}{{\'E}}1 {ù}{{\`u}}1 {û}{{\^u}}1 {î}{{\^i}}1 {ô}{{\^o}}1
           {«}{{\guillemotleft}}1 {»}{{\guillemotright}}1 {’}{{'}}1}
\lstset{style=emc}

% ---------------------------------------------------------------- titres
\usepackage{titlesec}
\titleformat{\chapter}[display]
  {\sffamily\bfseries}
  {\raggedright\normalsize\color{emcorange}\MakeUppercase{\chaptertitlename}~\thechapter}
  {2pt}
  {\Large\color{emcblue}\raggedright}
  [\vspace{2pt}{\color{emcborder}\titlerule[1.2pt]}]
\titlespacing*{\chapter}{0pt}{-22pt}{10pt}
\titleformat{\section}
  {\sffamily\bfseries\large\color{emcblue}}{\thesection}{0.6em}{}
\titlespacing*{\section}{0pt}{10pt}{4pt}
\titleformat{\subsection}
  {\sffamily\bfseries\normalsize\color{emcslate}}{\thesubsection}{0.6em}{}
\titlespacing*{\subsection}{0pt}{8pt}{2pt}

% ---------------------------------------------------------------- en-têtes
\usepackage{fancyhdr}
\pagestyle{fancy}
\fancyhf{}
\renewcommand{\headrulewidth}{0.4pt}
\renewcommand{\footrulewidth}{0pt}
\fancyhead[L]{\sffamily\scriptsize\color{emcgrey}EMC Helpline — Rapport de stage PFA}
\fancyhead[R]{\sffamily\scriptsize\color{emcgrey}\leftmark}
\fancyfoot[C]{\sffamily\small\color{emcblue}\thepage}
\fancypagestyle{plain}{\fancyhf{}%
  \renewcommand{\headrulewidth}{0pt}%
  \fancyfoot[C]{\sffamily\small\color{emcblue}\thepage}}
\renewcommand{\chaptermark}[1]{\markboth{\thechapter.~#1}{}}

% ---------------------------------------------------------------- glossaire
% \makenoidxglossaries : le glossaire est trié par LaTeX lui-même, sans
% makeglossaries ni Perl — le document compile partout avec xelatex + biber.
\usepackage[acronym,nomain,nopostdot,nogroupskip,nonumberlist,toc=false]{glossaries}
\makenoidxglossaries
\setacronymstyle{long-short}

% ---------------------------------------------------------------- biblio
\usepackage[backend=biber,style=numeric-comp,sorting=none,
            maxbibnames=4,giveninits=true,url=true,doi=false,isbn=false]{biblatex}
\addbibresource{references.bib}
\renewcommand*{\bibfont}{\small}

% ---------------------------------------------------------------- liens
\usepackage[hidelinks,breaklinks=true]{hyperref}
\hypersetup{
  colorlinks=true,linkcolor=emcblue,citecolor=emcorange,urlcolor=emcblue,
  pdftitle={EMC Helpline — Rapport de stage PFA},
  pdfauthor={Sayouba Ouedraogo},
  pdfsubject={Application mobile de signalement des cyberviolences},
  pdfkeywords={Flutter, Firebase, cyberviolence, protection de l'enfance, CMRPI},
  bookmarksnumbered=true,bookmarksopen=true}
\usepackage[french,capitalise,noabbrev]{cleveref}
\usepackage{url}
\urlstyle{sf}

% ---------------------------------------------------------------- utilitaires
\usepackage{xspace}
\usepackage{lastpage}
\newcommand{\emc}{EMC Helpline\xspace}

% Placeholder de capture d'écran.
%   \shot{<largeur>}{<chemin>}{<libellé court>}
% Si le fichier existe, il est inséré ; sinon un cadre pointillé indique
% exactement quel fichier déposer. Le rapport compile donc avant même que la
% moindre capture soit prise.
\newcommand{\shot}[4][2.05]{%
  \IfFileExists{#3}%
    {\includegraphics[width=#2]{#3}}%
    {\begin{tikzpicture}
       \node[draw=emcgrey!60,dashed,line width=0.7pt,fill=emcbg,
             rounded corners=3pt,minimum width=#2,
             minimum height={#1*#2},align=center,inner sep=3pt]
            {\sffamily\scriptsize\color{emcgrey}#4\\[3pt]
             \resizebox{\dimexpr#2-6pt\relax}{!}{\ttfamily\color{emcorange}\detokenize{#3}}};
     \end{tikzpicture}}}

% Capture au format téléphone (ratio 390:844)
\newcommand{\phoneshot}[3]{\shot[2.16]{#1}{#2}{#3}}
% Capture au format écran large (ratio 16:10)
\newcommand{\wideshot}[3]{\shot[0.62]{#1}{#2}{#3}}


% ------------------------------------------------------------------ SIGLES
\newacronym{cmrpi}{CMRPI}{Centre Marocain de Recherches Polytechniques et d'Innovation}
\newacronym{emc}{EMC}{Espace Maroc Cyberconfiance}
\newacronym{pfa}{PFA}{Projet de Fin d'Année}
\newacronym{ui}{UI}{Interface utilisateur (\emph{User Interface})}
\newacronym{ux}{UX}{Expérience utilisateur (\emph{User Experience})}
\newacronym{rtl}{RTL}{Écriture de droite à gauche (\emph{Right-To-Left})}
\newacronym{ttl}{TTL}{Durée de vie d'un enregistrement (\emph{Time To Live})}
\newacronym{sdk}{SDK}{Kit de développement (\emph{Software Development Kit})}
\newacronym{api}{API}{Interface de programmation (\emph{Application Programming Interface})}
\newacronym{fcm}{FCM}{Firebase Cloud Messaging}
\newacronym{ci}{CI}{Intégration continue (\emph{Continuous Integration})}
\newacronym{wcag}{WCAG}{\emph{Web Content Accessibility Guidelines}}
\newacronym{cndp}{CNDP}{Commission Nationale de contrôle de la protection des Données à caractère Personnel}
\newacronym{uid}{UID}{Identifiant unique d'utilisateur}
\newacronym{uml}{UML}{\emph{Unified Modeling Language}}

\begin{document}

% ============================================================================
%  PAGE DE GARDE
% ============================================================================
\begin{titlepage}
\thispagestyle{empty}
\centering

% --- Bandeau des logos -------------------------------------------------
% Déposer les deux fichiers dans figures/logos/ ; aucune autre modification
% n'est nécessaire, la page de garde les prend automatiquement.
\begin{minipage}[c]{0.42\textwidth}\centering
  \IfFileExists{figures/logos/ecole.png}
    {\includegraphics[height=2.4cm,keepaspectratio]{figures/logos/ecole.png}}
    {\begin{tikzpicture}
       \node[draw=emcgrey!50,dashed,fill=emcbg,rounded corners=2pt,
             minimum width=4.6cm,minimum height=2.4cm,align=center]
            {\sffamily\scriptsize\color{emcgrey}Logo de l'établissement\\[2pt]
             \tiny\ttfamily\color{emcorange}figures/logos/ecole.png};
     \end{tikzpicture}}
\end{minipage}\hfill
\begin{minipage}[c]{0.42\textwidth}\centering
  \IfFileExists{figures/logos/entreprise.png}
    {\includegraphics[height=2.4cm,keepaspectratio]{figures/logos/entreprise.png}}
    {\begin{tikzpicture}
       \node[draw=emcgrey!50,dashed,fill=emcbg,rounded corners=2pt,
             minimum width=4.6cm,minimum height=2.4cm,align=center]
            {\sffamily\scriptsize\color{emcgrey}Logo de l'organisme d'accueil\\[2pt]
             \tiny\ttfamily\color{emcorange}figures/logos/entreprise.png};
     \end{tikzpicture}}
\end{minipage}

\vspace{6mm}
{\sffamily\small\color{emcgrey}
 Université Cadi Ayyad — Faculté des Sciences et Techniques de Marrakech\\[1mm]
 Filière Ingénierie des Réseaux et Systèmes d'Information\par}

\vspace{4mm}
\noindent{\color{emcborder}\rule{\textwidth}{1pt}}
\vspace{8mm}

{\sffamily\bfseries\large\color{emcorange}RAPPORT DE PROJET DE FIN D'ANNÉE\par}
\vspace{3mm}
{\sffamily\small\color{emcgrey}Année universitaire 2025 — 2026\par}

\vspace{12mm}

{\sffamily\bfseries\LARGE\color{emcblue}
 Application mobile de signalement\\[2mm]
 des cyberviolences pour les enfants\par}

\vspace{5mm}
{\sffamily\large\color{emcslate}Ligne d'assistance \emc du \gls{cmrpi}\par}

\vspace{4mm}
{\itshape\color{emcgrey}\small
 Un prototype Flutter anonyme, trilingue et adossé à une architecture Firebase
 sans écriture directe depuis le client\par}

\vspace{14mm}
{\color{emcborder}\rule{0.5\textwidth}{0.6pt}}
\vspace{10mm}

\begin{minipage}[t]{0.46\textwidth}\raggedright
 {\sffamily\small\color{emcgrey}RÉALISÉ PAR\par}\vspace{2mm}
 {\sffamily\bfseries\large\color{emcslate}Sayouba Ouédraogo\par}
 {\small\color{emcgrey}Étudiant en 4\ieme{} année, IRISI\par}
\end{minipage}\hfill
\begin{minipage}[t]{0.46\textwidth}\raggedleft
 {\sffamily\small\color{emcgrey}ENCADRÉ PAR\par}\vspace{2mm}
 {\sffamily\bfseries\large\color{emcslate}Pr. Youssef Bentaleb\par}
 {\small\color{emcgrey}Encadrant principal, \gls{cmrpi}\par}\vspace{1.5mm}
 {\sffamily\bfseries\color{emcslate}Dr. Bouchra Assiss\par}
 {\small\color{emcgrey}Co-encadrante, \gls{cmrpi}\par}
\end{minipage}

\vfill
\noindent{\color{emcborder}\rule{\textwidth}{1pt}}\\[2mm]
{\sffamily\small\color{emcgrey}Stage effectué du 15 juillet au 31 août 2026\par}
\end{titlepage}

% ============================================================================
%  PAGES LIMINAIRES
% ============================================================================
\pagenumbering{roman}
\setcounter{page}{1}

\chapter*{Remerciements}
\addcontentsline{toc}{chapter}{Remerciements}
\markboth{Remerciements}{}

Je remercie le \gls{cmrpi} de m'avoir confié un sujet qui porte à conséquence :
une application dont se servent des enfants en difficulté n'autorise pas les
approximations d'usage, et cette exigence a structuré tout le stage.

Ma reconnaissance va au Professeur Youssef Bentaleb, encadrant principal, pour
la clarté du document de cadrage et pour des retours qui ont, à chaque jalon,
tranché des questions que je n'aurais pas su trancher seul — en particulier le
niveau de finition attendu des maquettes et le traitement des contenus
sensibles. Je remercie également la Docteure Bouchra Assiss pour son
co-encadrement.

Mes remerciements vont enfin au corps enseignant de la Faculté des Sciences et
Techniques de Marrakech, dont les enseignements en génie logiciel, en réseaux et
en sécurité ont trouvé dans ce projet un terrain d'application direct.

\vspace{4mm}
\begin{constat}
\small\sffamily
\textbf{Note de lecture.} Ce rapport rend compte d'un travail dont le périmètre
a dépassé celui du document de cadrage. Le \cref{tab:ecart} 
(\cref{chap:bilan}) confronte explicitement ce qui était demandé et ce qui a
été livré, afin que l'écart soit lisible et justifié plutôt que sous-entendu.
\end{constat}

\chapter*{Résumé}
\addcontentsline{toc}{chapter}{Résumé}
\markboth{Résumé}{}

Les enfants victimes de cyberviolences disposent au Maroc de canaux de
signalement conçus pour des adultes. Ce projet de fin d'année, mené au
\gls{cmrpi} du 15 juillet au 31 août 2026, a consisté à concevoir et développer
le prototype mobile de la ligne d'assistance \emc : une application Flutter par
laquelle un enfant dépose un signalement \textbf{sans compte et sans donner son
nom}, et reçoit en échange un numéro de référence — seul lien entre lui et son
dossier.

Le travail s'est déroulé en trois jalons de quinze jours : étude du public et
maquettage Figma, développement des écrans en Flutter, puis connexion au
back-end et finalisation. Le résultat dépasse le prototype à quatre écrans
initialement demandé : un parcours de signalement en onze étapes adapté à trois
tranches d'âge, une interface en français, arabe et anglais avec mise en miroir
complète en \gls{rtl}, un suivi de dossier par numéro de référence, une console
web pour l'équipe de traitement, et une politique de conservation de trente
jours glissants.

L'architecture retenue rompt avec le schéma Firebase usuel : \textbf{le client
n'écrit jamais dans la base}. Les règles Firestore refusent tout accès et chaque
opération passe par une fonction \emph{cloud} qui revalide le contrat de
données, applique une limitation de débit et exige une attestation d'intégrité
de l'application. Ce choix découle directement de l'absence de comptes
utilisateurs, elle-même dictée par la nature du public. Le numéro de référence
n'est stocké nulle part en clair : seule son empreinte SHA-256 sert de clé de
recherche.

L'ensemble est couvert par 290 tests automatisés exécutés en intégration
continue, dont un test de contrat qui compare les énumérations Dart au schéma
TypeScript et échoue dès que les deux divergent.

\vspace{3mm}
\noindent\textbf{Mots-clés} — cyberviolence, protection de l'enfance, Flutter,
Firebase, anonymat, accessibilité, architecture sans confiance client.

\vspace{6mm}
\noindent{\color{emcborder}\rule{\textwidth}{0.6pt}}
\vspace{4mm}

\begin{otherlanguage}{english}
\noindent\textbf{Abstract} — Children facing cyberviolence in Morocco are
offered reporting channels designed for adults. This end-of-year project,
carried out at \gls{cmrpi} between 15 July and 31 August 2026, delivers the
mobile prototype of the \emc line: a Flutter application through which a child
files a report \emph{without an account and without giving a name}, receiving a
reference number that is the only link between them and their case. The
delivered system exceeds the four-screen prototype initially scoped: an
eleven-step reporting flow adapted to three age groups, a French/Arabic/English
interface with full \gls{rtl} mirroring, reference-number case tracking, a web
console for the handling team, and a thirty-day rolling retention policy. The
architecture departs from the usual Firebase pattern: the client never writes to
the database. Security rules deny every direct access and each operation goes
through a cloud function that re-validates the data contract, rate-limits
requests and requires an app-integrity attestation. The reference number is
never stored in clear text — only its SHA-256 digest serves as a lookup key. The
whole is covered by 290 automated tests running in continuous integration.
\end{otherlanguage}

\cleardoublepage
\microtypesetup{protrusion=false}
\tableofcontents
% Les trois listes s'enchaînent sur la même page : \clearpage est neutralisé
% le temps de les composer, puis rétabli.
\begingroup
\let\originalclearpage\clearpage
\let\clearpage\relax
\vspace*{26pt}
\listoffigures
\vspace*{20pt}
\listoftables
\vspace*{20pt}
\printnoidxglossary[type=\acronymtype,title={Liste des acronymes},
                    toctitle={Liste des acronymes}]
\let\clearpage\originalclearpage
\endgroup
\microtypesetup{protrusion=true}

\cleardoublepage
\pagenumbering{arabic}
\setcounter{page}{1}

% ============================================================================
\chapter{Introduction générale}
\label{chap:intro}
% ============================================================================

Un enfant harcelé en ligne partage souvent son téléphone, parfois avec la
personne qu'il voudrait signaler. Il ne sait pas nommer ce qui lui arrive, il
craint qu'on prévienne ses parents, et il abandonne devant un formulaire qui lui
demande son identité. Les canaux de signalement existants au Maroc — le portail
national \texttt{e-vigilance}, les lignes téléphoniques — sont efficaces, mais
ils ont été pensés pour des adultes capables de qualifier un fait et d'assumer
une démarche administrative.

Le \gls{cmrpi}, à l'origine de la plateforme \gls{emc} et de la ligne
d'assistance \emc, a proposé de combler cet écart par une application mobile
dédiée. Le présent rapport rend compte du \gls{pfa} réalisé sur ce sujet du
15 juillet au 31 août 2026, en trois jalons de quinze jours.

\section{Problématique}

La difficulté du sujet n'est pas technique au sens habituel : construire quatre
écrans et écrire dans une base de données ne pose aucun problème d'ingénierie.
Elle tient à une contradiction qu'il faut résoudre, pas contourner.

\begin{constat}
Pour être utilisable par un enfant, l'application doit \textbf{ne rien demander} :
ni compte, ni nom, ni adresse électronique. Mais une application sans comptes est
une application sans identité, donc sans le mécanisme sur lequel repose toute la
sécurité usuelle d'une base de données \emph{cloud}. Comment garantir la
confidentialité de signalements d'enfants dans un système où personne ne
s'authentifie ?
\end{constat}

Cette question a commandé l'essentiel des choix d'architecture décrits au
\cref{chap:conception}, jusqu'à des conséquences que le projet a dû assumer,
comme le fait qu'un numéro de référence perdu soit un dossier définitivement
perdu.

\section{Objectifs du stage}

Le document de cadrage fixe un objectif volontairement mesuré : un prototype
mobile à quatre écrans, en français, sans intégration au back-end \gls{emc},
livré en trois jalons. Il s'agit d'un parcours d'apprentissage guidé plus que
d'un produit fini.

Le travail effectivement mené a suivi ce plan de jalons, mais l'a dépassé sur le
périmètre fonctionnel : le prototype livré est trilingue, comporte un parcours de
signalement en onze étapes, un suivi de dossier, une console de traitement et
une politique de conservation des données. Le \cref{chap:bilan} met cet écart en
regard des objectifs initiaux.

\section{Organisation du rapport}

Le \cref{chap:contexte} présente l'organisme d'accueil, l'analyse de l'existant
et la démarche. Le \cref{chap:besoins} spécifie les besoins, fonctionnels comme
non fonctionnels, et les contraintes réglementaires propres à un public mineur.
Le \cref{chap:conception} expose l'architecture, le modèle de données et la
conception de l'interface. Le \cref{chap:realisation} décrit la réalisation, du
parcours Flutter aux fonctions \emph{cloud}. Le \cref{chap:tests} rend compte de
la stratégie de test et de l'intégration continue. Le \cref{chap:bilan} dresse le
bilan et ouvre sur les extensions possibles.

% ============================================================================
\chapter{Contexte et cadre du projet}
\label{chap:contexte}
% ============================================================================

\section{Organisme d'accueil}

Le \gls{cmrpi} est une association scientifique marocaine qui opère sous l'égide
du ministère chargé de l'industrie et de l'économie numérique, avec l'appui du
Conseil de l'Europe. Il porte depuis 2013 la plateforme \gls{emc} et la campagne
nationale de prévention contre la cyberviolence et le cyberharcèlement
\cite{cyberconfiance,cmrpi}.

La ligne d'assistance \emc en constitue le volet opérationnel : elle accueille
les signalements de cyberviolence, propose un accompagnement juridique et
psychologique et intervient pour le retrait de contenus. Son extension
\emph{\emc for women}, lancée en décembre 2024, l'a ouverte aux femmes et jeunes
filles \cite{emcwomen}. Le présent projet en constitue le canal mobile destiné au
public le plus jeune.

\section{Analyse de l'existant}

Trois canaux coexistent avant ce projet, chacun avec une limite précise pour un
enfant.

\begin{table}[htbp]
\centering
\small
\caption{Canaux de signalement existants et limites pour un public mineur}
\label{tab:existant}
\begin{tabularx}{\textwidth}{@{}L{3.2cm} X L{4.3cm}@{}}
\toprule
\textbf{Canal} & \textbf{Fonctionnement} & \textbf{Limite pour un enfant} \\
\midrule
Portail \texttt{e-vigilance} & Formulaire web national de signalement des
cyberviolences \cite{evigilance} & Vocabulaire juridique, champs d'identité,
navigation dense sur mobile \\
Ligne téléphonique \emc & Contact direct avec un intervenant &
Suppose de parler, souvent impossible depuis un domicile partagé \\
Réseaux sociaux et signalement intégré aux plateformes & Retrait de contenu &
Aucun accompagnement, aucun suivi, aucune orientation locale \\
\bottomrule
\end{tabularx}
\end{table}

L'étude du portail \texttt{e-vigilance} a fourni le vocabulaire métier — types
d'incident, plateformes, niveaux d'urgence — repris tel quel dans les
énumérations de l'application, afin qu'un signalement mobile reste transposable
vers le système national.

\section{Objectifs et périmètre}

\begin{table}[htbp]
\centering
\small
\caption{Périmètre du prototype : ce qui est dans le stage et ce qui en est écarté}
\label{tab:perimetre}
\begin{tabularx}{\textwidth}{@{}L{1.1cm} X@{}}
\toprule
\textbf{Inclus} & Application mobile multilingue ; parcours de signalement adapté
aux tranches d'âge 5–8, 8–12 et 12–15 ans ; formulaire simple et rassurant ;
enregistrement du signalement dans Firebase \\
\midrule
\textbf{Écarté} & React Native ; intégration complète au back-end \gls{emc}
existant ; espace dédié aux parents et enseignants \\
\bottomrule
\end{tabularx}
\end{table}

\section{Démarche et planning}

Le stage a suivi les trois jalons de quinze jours du document de cadrage, chacun
clos par l'envoi d'un livrable à l'encadrant et par un retour écrit conditionnant
la suite ; la validation du jalon 2 a été prononcée sur la base d'une vidéo de
démonstration.

\begin{figure}[htbp]
\centering
\begin{tikzpicture}[font=\sffamily\scriptsize]
\begin{ganttchart}[
    hgrid=false, vgrid={*1{draw=emcborder}},
    x unit=0.247cm, y unit title=0.56cm, y unit chart=0.54cm,
    title height=1, title label font=\sffamily\scriptsize\color{emcgrey},
    bar height=0.44, bar top shift=0.28,
    bar/.append style={draw=none, fill=emcblue!75},
    bar label font=\sffamily\scriptsize\color{emcslate},
    group/.append style={draw=none, fill=emcorange},
    group label font=\sffamily\bfseries\scriptsize\color{emcslate},
    group left shift=0, group right shift=0, group top shift=0.5,
    group height=0.22, group peaks height=0,
    milestone/.append style={fill=emcred, draw=none},
    milestone label font=\sffamily\scriptsize\color{emcred},
    link/.style={-{Latex[length=1.2mm]}, emcgrey}
  ]{1}{48}
  \gantttitle{Juillet 2026}{17}
  \gantttitle{Août 2026}{31} \\
  \gantttitlelist{15,...,31}{1}
  \gantttitlelist{1,...,31}{1} \\
  \ganttgroup{Jalon 1 — Public et maquettes}{1}{17} \\
  \ganttbar{Ressources EMC, portail e-vigilance}{1}{5} \\
  \ganttbar{Besoins d'un jeune public}{4}{8} \\
  \ganttbar{Maquettage Figma (13 écrans)}{7}{17} \\
  \ganttgroup{Jalon 2 — Squelette Flutter}{18}{32} \\
  \ganttbar{Écrans et navigation}{18}{25} \\
  \ganttbar{Parcours en onze étapes, i18n, RTL}{24}{32} \\
  \ganttgroup{Jalon 3 — Back-end et finalisation}{33}{48} \\
  \ganttbar{Fonctions cloud, preuves, suivi}{33}{41} \\
  \ganttbar{Console, conservation, tests, CI}{40}{46} \\
  \ganttbar{Documentation et rapport}{44}{48} \\
  \ganttmilestone{L1}{17} \ganttmilestone{L2}{32} \ganttmilestone{L3}{48}
\end{ganttchart}
\end{tikzpicture}
\caption[Planning du stage en trois jalons]{Déroulement du stage. Les losanges
marquent les trois livrables (\textsf{L1} maquettes Figma, \textsf{L2}
application navigable, \textsf{L3} prototype fonctionnel documenté).}
\label{fig:gantt}
\end{figure}

% ============================================================================
\chapter{Analyse et spécification des besoins}
\label{chap:besoins}
% ============================================================================

\section{Le public cible et ce qu'il impose}

Le cadrage désigne trois tranches d'âge — 5 à 8 ans, 8 à 12 ans, 12 à 15 ans —
qui ne diffèrent pas seulement par le niveau de lecture. Elles diffèrent par ce
qu'un utilisateur est capable de faire seul, et l'analyse a converti ces
différences en règles de conception applicables.

\begin{table}[htbp]
\centering
\small
\caption{Des caractéristiques du public aux règles de conception}
\label{tab:public}
\begin{tabularx}{\textwidth}{@{}L{4.1cm} X@{}}
\toprule
\textbf{Caractéristique observée} & \textbf{Règle retenue} \\
\midrule
Lecture lente, vocabulaire limité & Une question par écran, phrases courtes,
choix illustrés plutôt que champs libres \\
Incapacité à qualifier juridiquement un fait & Types d'incident exprimés en
langage courant, énumérations fermées, option « autre » toujours ouverte \\
Téléphone partagé, parfois avec l'agresseur & Aucune donnée de signalement
écrite sur l'appareil ; notifications annoncées avant d'être demandées \\
Crainte d'être identifié ou dénoncé & Aucun compte, aucun nom réel demandé ;
pseudonyme facultatif \\
Danger immédiat possible & Numéro d'urgence affiché au-dessus de tout contenu
explicatif \\
Public plus âgé (12 ans et plus) autonome & Accès à un assistant conversationnel,
réservé à cette tranche \\
\bottomrule
\end{tabularx}
\end{table}

\section{Besoins fonctionnels}

\begin{figure}[htbp]
\centering
\begin{tikzpicture}[
  scale=0.88,transform shape,
  font=\sffamily\scriptsize,
  uc/.style={ellipse,draw=emcblue,fill=emcblue!6,line width=0.7pt,
             align=center,inner sep=2.2pt,minimum width=4.5cm,
             minimum height=1.02cm,text width=3.5cm},
  ucb/.style={uc,draw=emcorange,fill=emcorange!6},
  rel/.style={draw=emcgrey,line width=0.6pt},
  inc/.style={-{Latex[length=1.5mm]},draw=emcgrey,dashed,line width=0.6pt},
  actor/.pic={
    \draw[line width=0.8pt,emcslate] (0,0) circle (0.17);
    \draw[line width=0.8pt,emcslate] (0,-0.17) -- (0,-0.70)
      (-0.30,-0.34) -- (0.30,-0.34) (-0.26,-1.08) -- (0,-0.70) -- (0.26,-1.08);}]

% --- frontière du système
\node[draw=emcborder,line width=1pt,rounded corners=3pt,fill=emcbg,
      minimum width=10.1cm,minimum height=8.7cm] (sys) at (0,-0.25) {};
\node[anchor=north,font=\sffamily\bfseries\scriptsize\color{emcgrey},
      inner sep=4pt] at (sys.north) {Système \emc};

% --- cas d'utilisation du jeune public
\node[uc] (u1) at (-2.5, 3.05) {Déposer un signalement};
\node[uc] (u2) at (-2.5, 1.70) {Joindre des preuves};
\node[uc] (u3) at (-2.5, 0.35) {Suivre un dossier par numéro};
\node[uc] (u4) at (-2.5,-1.00) {Consulter les ressources};
\node[uc] (u5) at (-2.5,-2.35) {Contacter la ligne d'assistance};
\node[uc] (u6) at (-2.5,-3.70) {Dialoguer avec l'assistant (12 ans et plus)};

% --- cas d'utilisation de l'équipe
\node[ucb] (g1) at (2.6, 3.05) {Consulter un dossier};
\node[ucb] (g2) at (2.6, 1.70) {Changer le statut};
\node[ucb] (g3) at (2.6,-1.00) {Journaliser l'accès};

% --- acteur principal
\pic (a1) at (-6.9,0.6) {actor};
\node[align=center,font=\sffamily\scriptsize\color{emcslate}]
     at (-6.9,-0.95) {Enfant\\ ou jeune};
\foreach \n in {u1,u3,u4,u5,u6} \draw[rel] (-6.9,0.05) -- (\n.west);

% --- acteur secondaire
\pic (a2) at (6.9,2.4) {actor};
\node[align=center,font=\sffamily\scriptsize\color{emcslate}]
     at (6.9,0.85) {Agent \emc};
\draw[rel] (6.9,1.85) -- (g1.east);
\draw[rel] (6.9,1.85) -- (g2.east);

% --- relations
\draw[inc] (u1) -- node[right=1pt,font=\sffamily\tiny\color{emcgrey}]
      {\guillemotleft include\guillemotright} (u2);
\draw[inc] (g1.south) -- node[right=1pt,font=\sffamily\tiny\color{emcgrey}]
      {\guillemotleft include\guillemotright} (g2.north);
\draw[inc] (g2.south) -- (g3.north);
\end{tikzpicture}
\caption[Diagramme de cas d'utilisation]{Diagramme de cas d'utilisation \gls{uml}.
L'agent \gls{emc} n'agit pas depuis l'application mobile mais depuis la console
web (\cref{sec:console}) ; les deux s'adossent au même back-end, ce qui justifie
de les réunir dans une seule frontière.}
\label{fig:usecase}
\end{figure}

Le dépôt d'un signalement se décompose en onze étapes conditionnelles :
destinataire du signalement, tranche d'âge, genre, type d'incident, plateforme
concernée, preuves, besoin d'accompagnement, nature de l'accompagnement,
coordonnées, niveau d'urgence, puis récapitulatif. Trois règles de branchement
ne se devinent pas à la lecture et méritent d'être spécifiées explicitement.

\begin{table}[htbp]
\centering
\small
\caption{Règles de branchement du parcours de signalement}
\label{tab:branchement}
\begin{tabularx}{\textwidth}{@{}L{3.5cm} L{2.6cm} X@{}}
\toprule
\textbf{Réponse à « Veux-tu de l'aide ? »} & \textbf{Type d'aide} &
\textbf{Coordonnées demandées} \\
\midrule
Oui, un accompagnement & demandée & pseudonyme et téléphone, obligatoires \\
Je ne sais pas & étape sautée & étape sautée \\
Non & étape sautée & étape sautée \\
\bottomrule
\end{tabularx}
\end{table}

\begin{vigilance}
\small
L'étape des preuves ne peut pas être franchie à vide : il faut au moins une
capture, un lien, \textbf{ou} un récit d'au moins 120 caractères. Le récit tient
lieu de preuve parce que le cas le plus grave est fréquemment celui où
l'agresseur a effacé les traces. Rendre l'étape purement facultative aurait
donc écarté les signalements les plus urgents.
\end{vigilance}

\section{Besoins non fonctionnels}

\begin{table}[htbp]
\centering
\small
\caption{Exigences non fonctionnelles et moyen de vérification}
\label{tab:nfr}
\begin{tabularx}{\textwidth}{@{}L{0.9cm} L{3.5cm} X L{3.0cm}@{}}
\toprule
\textbf{Réf.} & \textbf{Exigence} & \textbf{Formulation} & \textbf{Vérification} \\
\midrule
NF1 & Anonymat & Aucun compte, aucun nom réel ; le pseudonyme et le téléphone
sont les seules données de contact & Revue de code, \cref{chap:conception} \\
NF2 & Confidentialité locale & Rien du contenu d'un signalement n'est écrit sur
l'appareil & Test automatisé dédié \\
NF3 & Confidentialité serveur & Le numéro de référence n'est jamais stocké en
clair & Revue du schéma, tests back-end \\
NF4 & Accessibilité & Contrastes conformes \gls{wcag}~2.1 niveau AA
\cite{wcag} ; cibles tactiles d'au moins 48~dp & Suite de tests
\texttt{accessibility\_test} \\
NF5 & Internationalisation & Français, arabe, anglais ; mise en miroir complète
en \gls{rtl} & Suite \texttt{rtl\_test} \\
NF6 & Résistance à l'abus & Limitation de débit et attestation d'intégrité de
l'application & Tests back-end \\
NF7 & Conservation & Trente jours glissants, preuves supprimées avec le dossier
& Déclencheur et politiques \gls{ttl} \\
NF8 & Robustesse d'envoi & Un réessai ne crée jamais deux dossiers & Clé
d'idempotence, tests \\
\bottomrule
\end{tabularx}
\end{table}

\section{Contraintes réglementaires et éthiques}

Trois cadres se superposent. La Convention relative aux droits de l'enfant
impose de protéger l'enfant contre toute forme de violence, y compris
numérique \cite{cide}. La loi 09-08 encadre le traitement des données à
caractère personnel au Maroc et soumet leur collecte à un principe de
minimisation \cite{loi0908}. La loi 103-13 réprime notamment la diffusion de
propos ou d'images d'une personne sans son consentement \cite{loi10313}, ce qui
concerne directement la catégorie d'incident « images intimes ».

\begin{decision}
\small
Ces trois cadres convergent vers une même conséquence technique : \textbf{ne
collecter que ce dont le traitement du dossier a besoin}. C'est la raison pour
laquelle le nom réel n'est demandé nulle part, et pour laquelle le caractère
anonyme d'un signalement n'est pas une case à cocher mais une conséquence — il
se déduit de l'absence de coordonnées.
\end{decision}

% ============================================================================
\chapter{Conception}
\label{chap:conception}
% ============================================================================

\section{La décision structurante : le client n'écrit pas dans la base}

Le schéma Firebase usuel place le \gls{sdk} client face à la base et confie la
sécurité aux règles Firestore, qui reposent sur une identité —
\lstinline[language=TypeScript]|request.auth.uid|. Or cette application n'a pas de
comptes, délibérément. Sans identité, la règle d'écriture dégénère en
\lstinline[language=TypeScript]|if true| et la clé d'\gls{api} est lisible dans le
paquet applicatif : n'importe qui peut écrire dans la collection, et une erreur
symétrique en lecture exposerait des signalements d'enfants. Le suivi par numéro
aggrave le problème — « lire un document si l'on connaît une chaîne » n'est pas
une règle Firestore, c'est un jeton porteur sans limitation de débit.

\begin{decision}
\small
Les règles Firestore et Cloud Storage refusent tout à tout le monde
(\lstinline[language=TypeScript]|allow read, write: if false|) et chaque
opération passe par une fonction \emph{cloud} appelable. Conséquence pratique :
aucune dépendance \texttt{cloud\_firestore} côté Flutter — le client ne connaît
que des appels de fonctions.
\end{decision}

\begin{figure}[htbp]
\centering
\begin{tikzpicture}[
  scale=0.92,transform shape,
  font=\sffamily\scriptsize,
  layer/.style={draw=emcborder,line width=0.9pt,rounded corners=3pt,fill=white},
  blk/.style={draw=emcblue!45,line width=0.6pt,rounded corners=2pt,
              fill=emcblue!7,align=center,inner sep=3pt,
              minimum height=0.78cm,text width=2.25cm},
  fn/.style={blk,draw=emcorange!55,fill=emcorange!8},
  db/.style={blk,draw=emcgreen!55,fill=emcgreen!7},
  lbl/.style={font=\sffamily\bfseries\scriptsize\color{emcgrey},anchor=west},
  flow/.style={-{Latex[length=1.6mm]},draw=emcgrey,line width=0.8pt}]

% ---- couche client
\node[layer,minimum width=12.6cm,minimum height=1.75cm] (cl) at (0,3.55) {};
\node[lbl] at ([xshift=5pt,yshift=-7pt]cl.north west) {Application Flutter (Android / iOS)};
\node[blk] (v) at (-4.4,3.25) {\texttt{views/}\\ 11 étapes, 4 onglets};
\node[blk] (p) at (-1.47,3.25) {\texttt{providers/}\\ état et validation};
\node[blk] (m) at ( 1.47,3.25) {\texttt{models/}\\ énumérations métier};
\node[blk] (c) at ( 4.4,3.25) {\texttt{core/}\\ i18n, envoi, stockage};

% ---- frontière de confiance
\node[layer,minimum width=12.6cm,minimum height=1.30cm,fill=emcbg] (fr) at (0,1.55) {};
\node[lbl] at ([xshift=5pt,yshift=-7pt]fr.north west) {Frontière de confiance};
\node[blk,text width=3.1cm,draw=emcred!45,fill=emcred!6] (ac) at (-2.4,1.25)
     {App Check\\ attestation d'intégrité};
\node[blk,text width=3.1cm,draw=emcred!45,fill=emcred!6] (an) at ( 2.0,1.25)
     {Auth anonyme\\ clé de comptage};

% ---- couche fonctions
\node[layer,minimum width=12.6cm,minimum height=1.75cm] (fc) at (0,-0.7) {};
\node[lbl] at ([xshift=5pt,yshift=-7pt]fc.north west) {Cloud Functions 2\ieme{} génération (TypeScript, Node 22)};
\node[fn] (f1) at (-4.4,-1.0) {\texttt{submitReport}};
\node[fn] (f2) at (-1.47,-1.0) {\texttt{trackReport}};
\node[fn] (f3) at ( 1.47,-1.0) {\texttt{requestEvidence-}\\\texttt{UploadUrl}};
\node[fn] (f4) at ( 4.4,-1.0) {\texttt{console*}\\ \texttt{retention}};

% ---- couche données
\node[layer,minimum width=12.6cm,minimum height=1.75cm] (dl) at (0,-2.95) {};
\node[lbl] at ([xshift=5pt,yshift=-7pt]dl.north west) {Données — accès direct refusé par les règles};
\node[db] (d1) at (-4.4,-3.25) {\texttt{reports}\\ \texttt{referenceIndex}};
\node[db] (d2) at (-1.47,-3.25) {\texttt{idempotency}\\ \texttt{rateLimits}};
\node[db] (d3) at ( 1.47,-3.25) {\texttt{auditLog}};
\node[db] (d4) at ( 4.4,-3.25) {Cloud Storage\\ \texttt{evidence/}};

% ---- second client : la console de l'équipe
\node[blk,text width=2.3cm,draw=emcslate!45,fill=emcslate!6,minimum height=1.0cm]
     (cs) at (7.35,1.9) {Console équipe\\ (web)};
\node[font=\sffamily\tiny\color{emcgrey},align=center,text width=2.2cm]
     at (7.35,0.75) {Auth par courriel\\ + rôle \texttt{agent}};

\draw[flow] (cl.south) -- (fr.north);
\draw[flow] (fr.south) -- (fc.north);
\draw[flow] (fc.south) -- (dl.north);
\draw[flow] (cs.south) -- (7.35,-0.7) -- (fc.east);
\end{tikzpicture}
\caption[Architecture applicative]{Architecture en couches. Toute flèche
entrante dans la couche de données passe obligatoirement par une fonction
\emph{cloud} : il n'existe aucun chemin direct du client vers Firestore.}
\label{fig:archi}
\end{figure}

\section{Modèle de données}

\begin{figure}[htbp]
\centering
\begin{tikzpicture}[
  scale=0.94,transform shape,
  font=\sffamily\scriptsize,
  boxbase/.style={rectangle split,draw,line width=0.7pt,fill=white,
                  align=left,inner sep=3.2pt,
                  every one node part/.style={font=\sffamily\bfseries\scriptsize,align=center},
                  every two node part/.style={font=\ttfamily\tiny},
                  every three node part/.style={font=\ttfamily\tiny}},
  cls/.style={boxbase,rectangle split parts=3,draw=emcblue,
              rectangle split part fill={emcblue!12,white,white}},
  enum/.style={boxbase,rectangle split parts=2,draw=emcorange,
              rectangle split part fill={emcorange!12,white},text width=3.35cm},
  iface/.style={boxbase,rectangle split parts=2,draw=emcgreen,
              rectangle split part fill={emcgreen!12,white},text width=3.5cm},
  impl/.style={boxbase,rectangle split parts=3,draw=emcslate,
              rectangle split part fill={emcslate!12,white,white},text width=3.5cm},
  assoc/.style={draw=emcgrey,line width=0.6pt},
  real/.style={dashed,draw=emcgrey,line width=0.6pt,-{Triangle[open,length=2.2mm,width=2.2mm]}},
  dep/.style={draw=emcgrey,line width=0.6pt,dashed,-{Latex[length=1.5mm]}}]

\node[cls,text width=3.5cm,anchor=north] (rp) at (-5.7,3.1) {%
  ReportProvider
  \nodepart{two}
  -- step : int\\
  -- report : ReportModel\\
  -- idempotencyKey : String
  \nodepart{three}
  + next() / back()\\
  + isStepValid() : bool\\
  + submit() : SubmissionOutcome};

\node[cls,text width=4.15cm,anchor=north] (rm) at (0.35,3.1) {%
  ReportModel
  \nodepart{two}
  + whoFor : WhoFor\\
  + ageGroup : AgeGroup\\
  + gender : Gender\\
  + incidentType : IncidentType\\
  + platform : ReportPlatform\\
  + description : String?\\
  + evidencePaths : List<String>\\
  + pseudo, contactPhone : String?\\
  + urgencyLevel : UrgencyLevel
  \nodepart{three}
  + isAnonymous : bool\\
  + toPayload() : Map};

\node[enum,anchor=north] (e1) at (6.05,3.1) {%
  \guillemotleft enumeration\guillemotright\\ IncidentType
  \nodepart{two}
  hateSpeech, discrimination,\\
  defamation, identityTheft,\\
  intimateImages, threat, other};

\node[enum,anchor=north] (e2) at (6.05,0.45) {%
  \guillemotleft enumeration\guillemotright\\ AgeGroup
  \nodepart{two}
  child, teen, adult, undisclosed\\
  + isChatbotEligible : bool};

\node[enum,anchor=north] (e3) at (6.05,-1.75) {%
  \guillemotleft enumeration\guillemotright\\ AssistanceNeed
  \nodepart{two}
  wanted, none, unsure};

\node[iface,anchor=north] (if) at (-5.7,-0.45) {%
  \guillemotleft interface\guillemotright\\ ReportSubmitter
  \nodepart{two}
  + submit(payload, key)\\
  + track(code)};

\node[impl,anchor=north] (fb) at (-5.7,-2.55) {%
  FirebaseBackend
  \nodepart{two}
  -- functions : CloudFunctions
  \nodepart{three}
  + submit(...) \quad + track(...)};

\draw[assoc] (rp.east|-rm.west) ++(0,0) coordinate (tmp);
\draw[assoc] ([yshift=-6mm]rp.north east) -- node[above,font=\sffamily\tiny\color{emcgrey}]{1}
             ([yshift=-6mm]rm.north west);
\draw[assoc] ([yshift=-8mm]rm.north east) -- ([yshift=-8mm]e1.north west);
\draw[assoc] (rm.east) -- (e2.west);
\draw[assoc] ([yshift=6mm]rm.south east) -- (e3.west);
\draw[real] (fb.north) -- (if.south);
\draw[dep] (rp.south) -- (if.north);
\end{tikzpicture}
\caption[Diagramme de classes du domaine]{Diagramme de classes \gls{uml}
(extrait). Le provider ne connaît que l'interface \texttt{ReportSubmitter} :
l'implémentation Firebase est injectée, ce qui permet de tester tout le parcours
sans réseau.}
\label{fig:classes}
\end{figure}

Côté serveur, cinq collections se répartissent les responsabilités
(\cref{tab:collections}), la séparation du dossier et de sa clé d'accès en étant
le point central.

\begin{table}[htbp]
\centering
\small
\caption{Collections Firestore et rôle de chacune}
\label{tab:collections}
\begin{tabularx}{\textwidth}{@{}L{3.35cm} X L{2.15cm}@{}}
\toprule
\textbf{Collection} & \textbf{Contenu} & \textbf{Expiration} \\
\midrule
\texttt{reports/\{id\}} & Le signalement, son statut et les chemins des preuves &
30 j glissants \\
\texttt{referenceIndex/}\linebreak\texttt{\{sha256\}} & Empreinte du numéro de référence
$\rightarrow$ identifiant du dossier & en phase \\
\texttt{idempotency/\{clé\}} & Numéro déjà attribué, rendu tel quel si la même
clé revient & 30 j \\
\texttt{rateLimits/\{clé\}} & Compteurs de la limitation de débit & fenêtre
courte \\
\texttt{auditLog/\{id\}} & Qui a lu ou modifié quel dossier, sans aucun
contenu & survit au dossier \\
\bottomrule
\end{tabularx}
\end{table}

\begin{decision}
\small
Le numéro de référence est un secret porteur : il ouvre le dossier. Son
empreinte SHA-256 sert donc d'\textbf{identifiant de document} dans
\texttt{referenceIndex}, ce qui donne trois propriétés d'un coup : une contrainte
d'unicité réelle, une lecture par clé exacte sans index composite, et la
séparation du dossier et de sa clé d'accès. Aucun sel n'est ajouté — la valeur
hachée porte 60 bits d'aléa, et un sel par enregistrement rendrait la recherche
par code impossible, ce qui est précisément son objet.
\end{decision}

\begin{vigilance}
\small
La contrepartie est assumée et affichée dans l'application : \textbf{un code
perdu est un dossier perdu}. C'est le prix de l'anonymat, et il vaut mieux
l'énoncer à l'utilisateur que le découvrir au moment du suivi.
\end{vigilance}

\section{Déroulement d'un dépôt}

\begin{figure}[htbp]
\centering
\begin{tikzpicture}[
  scale=0.92,transform shape,
  font=\sffamily\scriptsize,
  ll/.style={draw=emcborder,line width=0.7pt,dash pattern=on 2pt off 2pt},
  head/.style={draw=emcblue,line width=0.7pt,fill=emcblue!8,rounded corners=2pt,
               align=center,inner sep=3pt,text width=2.15cm,minimum height=0.62cm},
  act/.style={draw=emcblue!60,fill=emcblue!18,line width=0.5pt},
  msg/.style={-{Latex[length=1.5mm]},draw=emcslate,line width=0.65pt},
  ret/.style={-{Latex[length=1.5mm]},draw=emcgrey,line width=0.6pt,dashed},
  note/.style={font=\sffamily\tiny\color{emcgrey},align=left}]

\def\ya{0}      % top of lifelines
\def\yb{-6.7}   % bottom
\foreach \x/\n [count=\i] in {-6.1/App Flutter, -2.6/App Check, 0.9/{\texttt{submitReport}},
                              4.2/Firestore, 7.0/Cloud Storage} {
  \node[head] (h\i) at (\x,\ya) {\n};
  \draw[ll] (\x,\ya-0.34) -- (\x,\yb);
}

% barres d'activation
\fill[act] (-6.22,-0.85) rectangle (-5.98,-6.35);
\fill[act] ( 0.78,-2.8) rectangle ( 1.02,-5.55);

\draw[msg] (-6.1,-1.1) -- (-6.1,-1.1);
\node[note,anchor=west] at (-5.9,-1.0) {1.~Onze étapes remplies, validation locale à chaque écran};

\draw[msg] (-6.0,-1.7) -- node[above,note]{2.~URL signée} (6.9,-1.7);
\draw[ret] (7.0,-2.1) -- node[above,note]{3.~téléversement direct de la capture} (-6.0,-2.1);

\draw[msg] (-6.0,-2.8) -- node[above,note]{4.~jeton d'intégrité} (-2.7,-2.8);
\draw[ret] (-2.6,-3.2) -- (-6.0,-3.2);

\draw[msg] (-6.0,-3.75) -- node[above,note]{5.~\texttt{submitReport(payload, clé d'idempotence)}} (0.78,-3.75);
\node[note,anchor=west] at (1.2,-4.2)
     {6.~revalidation du schéma, limitation de débit,\\ \phantom{6.~}génération d'un code de 60 bits};

\draw[msg] (1.02,-4.9) -- node[above,note]{7.~transaction : dossier + index + idempotence} (4.1,-4.9);
\draw[ret] (4.2,-5.3) -- (1.02,-5.3);

\draw[ret] (0.78,-5.8) -- node[above,note]{8.~numéro de référence} (-6.0,-5.8);
\node[note,anchor=west] at (-5.9,-6.25) {9.~Affichage du numéro ; rien n'est écrit sur l'appareil};

\node[draw=emcorange!50,fill=emcorange!6,rounded corners=2pt,line width=0.6pt,
      text width=3.5cm,align=left,font=\sffamily\tiny\color{emcslate},inner sep=3pt]
     at (5.9,-6.25)
     {Si la même clé d'idempotence revient, l'étape 7 est court-circuitée et le
      même numéro est rendu : un réessai ne crée jamais deux dossiers.};
\end{tikzpicture}
\caption[Diagramme de séquence du dépôt]{Diagramme de séquence \gls{uml} du dépôt.
Les preuves ne passent pas par la fonction : le client les téléverse directement
vers Cloud Storage au moyen d'une URL signée, de durée de vie limitée.}
\label{fig:sequence}
\end{figure}

\section{Conception de l'interface}

Le jalon 1 a produit treize écrans sous Figma, accompagnés de leurs variantes de
sélection — l'état d'un écran après qu'un choix a été fait —, afin que la
maquette exprime le comportement et pas seulement la mise en page.

\begin{figure}[htbp]
\centering
\begin{subfigure}[b]{0.268\textwidth}\centering
  \phoneshot{\linewidth}{figures/screens/figma-01-accueil.png}{Maquette — accueil}
  \caption{Accueil}
\end{subfigure}\hfill
\begin{subfigure}[b]{0.268\textwidth}\centering
  \phoneshot{\linewidth}{figures/screens/figma-05-incident.png}{Maquette — type d'incident}
  \caption{Type d'incident}
\end{subfigure}\hfill
\begin{subfigure}[b]{0.268\textwidth}\centering
  \phoneshot{\linewidth}{figures/screens/figma-10-recap.png}{Maquette — récapitulatif}
  \caption{Récapitulatif}
\end{subfigure}
\caption[Maquettes Figma du parcours]{Trois des treize écrans maquettés au jalon 1 ;
la charte a été reprise telle quelle dans le code.}
\label{fig:figma}
\end{figure}

L'accessibilité a corrigé la charte issue des maquettes : l'orange d'origine
\texttt{\#EA580C} n'atteignait que 3{,}56:1 de contraste sur blanc, en deçà des
4{,}5:1 exigés par le niveau AA pour du texte courant \cite{wcag}. Il a été
assombri en \texttt{\#C2410C} (5{,}18:1) pour le texte et les boutons, la teinte
vive restant réservée aux éléments décoratifs, qui n'exigent que 3:1 ; le rouge
d'alerte a subi le même traitement.

% ============================================================================
\chapter{Réalisation}
\label{chap:realisation}
% ============================================================================

\section{Organisation du dépôt}

Le projet réunit trois parties dans un dépôt unique.

\begin{table}[htbp]
\centering
\small
\caption{Volumétrie du code livré (hors code généré et dépendances)}
\label{tab:volumetrie}
\begin{tabularx}{\textwidth}{@{}L{3.1cm} L{2.55cm} L{1.7cm} X@{}}
\toprule
\textbf{Partie} & \textbf{Emplacement} & \textbf{Lignes} & \textbf{Technologie} \\
\midrule
Application mobile & \texttt{lib/} & 14 500 & Flutter 3.44.6 / Dart 3.12 \\
Fonctions \emph{cloud} & \texttt{functions/src/} & 1 640 & TypeScript, Node 22 \\
Tests applicatifs & \texttt{test/} & 3 390 & \texttt{flutter\_test} \\
Console équipe & \texttt{console/} & --- & Page statique, Firebase Hosting \\
\bottomrule
\end{tabularx}
\end{table}

La séparation est stricte : \texttt{views/} ne porte que de la présentation,
\texttt{providers/} l'état du parcours et sa validation, \texttt{models/} les
énumérations métier, \texttt{core/} l'accès au back-end, la localisation et les
préférences.

\section{Le parcours de signalement}

Les onze étapes partagent une ossature commune — \texttt{StepLayout} — qui fixe la
position du titre, de la barre de progression et des boutons. Le provider expose
une validation par étape : le bouton « continuer » reste inactif tant que la
réponse manque, plutôt que d'afficher une erreur après coup.

\begin{figure}[htbp]
\centering
\begin{subfigure}[b]{0.213\textwidth}\centering
  \phoneshot{\linewidth}{figures/screens/app-01-accueil.png}{App — accueil}
  \caption{Accueil}\end{subfigure}\hfill
\begin{subfigure}[b]{0.213\textwidth}\centering
  \phoneshot{\linewidth}{figures/screens/app-02-age.png}{App — tranche d'âge}
  \caption{Tranche d'âge}\end{subfigure}\hfill
\begin{subfigure}[b]{0.213\textwidth}\centering
  \phoneshot{\linewidth}{figures/screens/app-03-preuves.png}{App — preuves}
  \caption{Preuves}\end{subfigure}\hfill
\begin{subfigure}[b]{0.213\textwidth}\centering
  \phoneshot{\linewidth}{figures/screens/app-04-recap.png}{App — récapitulatif}
  \caption{Récapitulatif}\end{subfigure}
\caption[Parcours de signalement dans l'application]{Quatre écrans du parcours. Le
récapitulatif signale qu'un récit a été écrit sans le réafficher : sur un
téléphone partagé, le relire à l'écran est un risque.}
\label{fig:parcours}
\end{figure}

\begin{vigilance}
\small
L'écran de premier lancement \textbf{ne demande aucune permission}. Android
cesse d'afficher la boîte de dialogue des notifications après deux refus, et un
« non » exprimé avant même de savoir à quoi sert l'application la condamne
définitivement. Le choix retenu : annoncer d'avance, demander au moment où il y
a enfin quelque chose à notifier. Le numéro d'urgence est placé au-dessus du
texte de présentation — qui est en danger immédiat n'a pas à lire une
introduction.
\end{vigilance}

\section{Trois langues et la mise en miroir}

L'interface est déclinée en français, arabe et anglais au moyen de trois fichiers
ARB et de la génération \texttt{AppLocalizations} de Flutter. En arabe, ce n'est
pas seulement le texte qui change : toute la mise en page est mise en miroir. Les
éléments que Flutter ne retourne pas seuls — icônes directionnelles, décalages
calculés, barres de progression — ont demandé un traitement explicite, couvert
par une suite de tests dédiée \cite{flutterdocs}.

\begin{decision}
\small
Les libellés d'affichage ne vivent pas dans les énumérations. Celles-ci ne
portent aucun texte : \texttt{IncidentType.threat} est un symbole, et son libellé
est résolu dans \texttt{report\_enum\_labels.dart} à partir des fichiers ARB.
Traduire une étiquette ne peut donc pas modifier le branchement du parcours — ce
qui était le cas tant que la comparaison portait sur des chaînes françaises.
\end{decision}

\section{Le contrat de données, tenu des deux côtés}

Les noms des membres d'énumération voyagent tels quels du client Dart jusqu'au
schéma TypeScript ; une divergence casserait l'envoi silencieusement, en
production seulement. Un test de contrat lit le fichier TypeScript depuis la
suite Dart et échoue dès que les deux listes ne coïncident plus.

\begin{lstlisting}[language=Dart,caption={Le test de contrat compare les
énumérations Dart au schéma serveur},label={lst:contrat}]
// test/backend_contract_test.dart
test('les valeurs d incident correspondent au schema TypeScript', () {
  final schema = File('functions/src/schema.ts').readAsStringSync();
  final serveur = _extraireUnion(schema, 'incidentType');
  final client  = IncidentType.values.map((e) => e.name).toSet();
  expect(client, equals(serveur));   // echoue des que les listes divergent
});
\end{lstlisting}

\section{Dépôt, idempotence et limitation de débit}

La fonction \texttt{submitReport} ne fait pas confiance au client : elle revalide
la charge utile, applique la limitation de débit, puis écrit le dossier, son
index et la trace d'idempotence en une seule transaction.

\begin{lstlisting}[language=TypeScript,caption={Ecriture transactionnelle du
dossier et de son index (extrait simplifié)},label={lst:submit}]
// functions/src/submitReport.ts
const code = genererCodeReference();            // 60 bits d'aleatoire
const empreinte = sha256(normaliser(code));     // seule valeur persistee

await db.runTransaction(async (tx) => {
  const dejaVu = await tx.get(db.doc(`idempotency/${cle}`));
  if (dejaVu.exists) return dejaVu.data().referenceCode;   // reessai : meme code

  const dossier = db.collection('reports').doc();
  tx.set(dossier, { ...charge, status: 'received',
                    createdAt: now, expiresAt: dans30Jours });
  tx.set(db.doc(`referenceIndex/${empreinte}`),
         { reportId: dossier.id, createdAt: now, expiresAt: dans30Jours });
  tx.set(db.doc(`idempotency/${cle}`),
         { reportId: dossier.id, referenceCode: code, expiresAt: dans30Jours });
});
\end{lstlisting}

\begin{figure}[htbp]
\centering
\begin{subfigure}[b]{0.213\textwidth}\centering
  \phoneshot{\linewidth}{figures/screens/app-05-envoi.png}{App — envoi en cours}
  \caption{Envoi}\end{subfigure}\hfill
\begin{subfigure}[b]{0.213\textwidth}\centering
  \phoneshot{\linewidth}{figures/screens/app-06-succes.png}{App — numéro de référence}
  \caption{Numéro rendu}\end{subfigure}\hfill
\begin{subfigure}[b]{0.213\textwidth}\centering
  \phoneshot{\linewidth}{figures/screens/app-07-suivi.png}{App — suivi par numéro}
  \caption{Suivi}\end{subfigure}\hfill
\begin{subfigure}[b]{0.213\textwidth}\centering
  \phoneshot{\linewidth}{figures/screens/app-08-arabe.png}{App — accueil en arabe (RTL)}
  \caption{Arabe, \gls{rtl}}\end{subfigure}
\caption[Envoi, suivi et version arabe]{Fin du parcours et suivi de dossier ; le
quatrième écran montre l'interface en arabe, entièrement mise en miroir.}
\label{fig:suivi}
\end{figure}

\section{Les preuves}
\label{sec:preuves}

Les captures ne transitent jamais par la fonction de dépôt : le client demande
une URL de téléversement signée, envoie le fichier directement à Cloud Storage,
puis ne transmet que le chemin obtenu. La fonction reste ainsi légère quel que
soit le poids des pièces, et un fichier téléversé pour un signalement jamais
envoyé finit supprimé par une règle de cycle de vie.

\begin{vigilance}
\small
La politique \gls{ttl} de Firestore ignore Cloud Storage. Sans le déclencheur
\texttt{onReportDeleted}, l'expiration effacerait le dossier en laissant les
captures — des images de l'agression d'un enfant — dans le \emph{bucket}, pour
toujours, détachées de l'enregistrement qui expliquait leur présence. Vue depuis
Firestore, cette panne est invisible : tout y paraît conforme.
\end{vigilance}

\section{La console de l'équipe et la conservation}
\label{sec:console}

Les agents \gls{emc} travaillent depuis une application web distincte, authentifiée
par courriel et autorisée par un \emph{custom claim}
\lstinline[language=TypeScript]|role: agent|. Elle liste les dossiers, fait évoluer
leur statut — \texttt{received}, \texttt{inReview}, \texttt{contacted},
\texttt{closed} — et affiche le compte à rebours de conservation, en rouge sous
sept jours. Chaque lecture et chaque changement d'état alimentent un journal
d'audit, dépourvu de tout contenu de signalement et qui survit au dossier.

\begin{decision}
\small
La conservation est de \textbf{trente jours glissants} : le dépôt la fixe, et
chaque changement de statut la repousse. Trente jours fixes depuis la création
supprimeraient un dossier pendant son instruction — le seul comportement qu'une
politique de conservation ne doit jamais avoir.
\end{decision}

\noindent La console est reproduite en \cref{fig:validation}.

% ============================================================================
\chapter{Tests, intégration continue et validation}
\label{chap:tests}
% ============================================================================

\section{Stratégie de test}

L'application vise des enfants en détresse, et une régression silencieuse y coûte
plus cher qu'ailleurs. La couverture a donc été bâtie autour des propriétés qu'il
serait grave de perdre, plutôt que d'un pourcentage de lignes. Au total,
\textbf{187 tests côté application et 103 côté back-end}, soit 290 tests
automatisés ; les suites serveur s'exécutent contre les émulateurs Firebase, ce
qui permet de vérifier transactions, règles de sécurité et idempotence sans
compte de facturation ni donnée réelle.

\begin{table}[htbp]
\centering
\small
\caption{Suites de tests et propriété vérifiée}
\label{tab:tests}
\begin{tabularx}{\textwidth}{@{}L{4.6cm} X@{}}
\toprule
\textbf{Suite} & \textbf{Propriété garantie} \\
\midrule
\texttt{report\_provider\_test} & Le branchement des onze étapes, sauts conditionnels compris \\
\texttt{wizard\_steps\_test} & Une étape refuse d'être franchie tant que sa réponse manque \\
\texttt{backend\_contract\_test} & Les énumérations Dart et le schéma TypeScript
ne divergent pas (\cref{lst:contrat}) \\
\texttt{rtl\_test} & La mise en miroir en arabe atteint ce que Flutter ne retourne pas seul \\
\texttt{accessibility\_test} & Contrastes et tailles de cibles tactiles \\
\texttt{settings\_store\_test} & Aucune donnée de signalement n'atteint le disque \\
\texttt{submission\_error\_test} & Un réessai après échec ne crée pas un second
dossier \\
\texttt{firestore.rules.test} & Toute écriture directe depuis un client est
refusée \\
\bottomrule
\end{tabularx}
\end{table}

\section{Validation de bout en bout}

Le parcours complet — dépôt, téléversement d'une capture, attribution d'un numéro,
suivi, traitement depuis la console — a été rejoué en local sur la pile
d'émulateurs : le signalement de test apparaît dans \texttt{reports}, son empreinte
dans \texttt{referenceIndex}, et la console le retrouve avec son compte à rebours.

\begin{vigilance}
\small
Deux comportements ne sont pas observables en local et restent à valider après
déploiement : la signature des URL de téléversement, que les émulateurs
n'effectuent pas, et les notifications \gls{fcm}, dont le code est en place mais
qui exigent un projet réel. Ces limites sont documentées plutôt que passées sous
silence.
\end{vigilance}

\section{Intégration continue}

Deux tâches indépendantes s'exécutent à chaque \emph{push} : analyse statique,
formatage et tests Flutter d'une part ; contrôle de types TypeScript et suites
back-end contre les émulateurs d'autre part.

\begin{decision}
\small
La version de Flutter est \textbf{épinglée}, en intégration continue comme en
local. Le formateur de Dart change de règles d'une version à l'autre : sans
épingle, un fichier correct devient incorrect le jour où le canal stable avance,
et la chaîne échoue sur un commit qui n'a rien touché. Monter de version devient
ainsi un acte délibéré, qui reformate ce qu'il faut.
\end{decision}

\begin{figure}[htbp]
\centering
\begin{subfigure}[b]{0.455\textwidth}\centering
  \wideshot{\linewidth}{figures/screens/console-dossier.png}{Console — dossier et journal d'audit}
  \caption{Console de l'équipe}\end{subfigure}\hfill
\begin{subfigure}[b]{0.455\textwidth}\centering
  \wideshot{\linewidth}{figures/screens/ci-github-actions.png}{Intégration continue — deux tâches au vert}
  \caption{Intégration continue}\end{subfigure}
\caption[Console de traitement et intégration continue]{La console de l'équipe
\gls{emc} et les deux tâches d'intégration continue exécutées sur chaque \emph{push}.}
\label{fig:validation}
\end{figure}

% ============================================================================
\chapter{Bilan et perspectives}
\label{chap:bilan}
% ============================================================================

\section{Ce qui était demandé, ce qui a été livré}

\begin{table}[htbp]
\centering
\small
\caption{Objectifs du document de cadrage confrontés au livrable}
\label{tab:ecart}
\begin{tabularx}{\textwidth}{@{}L{4.0cm} L{4.4cm} X@{}}
\toprule
\textbf{Attendu (cadrage)} & \textbf{Livré} & \textbf{Écart} \\
\midrule
Prototype à 4 écrans & 4 onglets, 11 étapes de signalement, 7 écrans annexes &
Parcours complet plutôt qu'ossature \\
Français d'abord, arabe reporté & Français, arabe, anglais, \gls{rtl} complet &
Extension traitée pendant le stage \\
Adaptation aux tranches d'âge & 4 groupes d'âge, assistant réservé aux 12 ans et
plus & Conforme \\
Enregistrement dans Firebase & Fonctions \emph{cloud}, écriture directe interdite
& Architecture durcie \\
Documentation du prototype & README, plan back-end, 290 tests, \gls{ci} &
Conforme et étendu \\
Back-end \gls{emc} écarté & Écarté & Conforme \\
Espace parents et enseignants écarté & Écarté & Conforme \\
\bottomrule
\end{tabularx}
\end{table}

Trois éléments s'ajoutent hors cadrage : le suivi de dossier par numéro de
référence, la console de traitement pour l'équipe, et la politique de
conservation à trente jours glissants. Chacun répond à une question qui se pose
dès qu'un signalement est réellement enregistré : comment l'auteur en a-t-il des
nouvelles, qui le traite, et combien de temps le conserve-t-on.

\section{Limites connues}

Elles sont documentées dans le dépôt et reprises ici sans atténuation.

\begin{itemize}
\item \textbf{L'assistant conversationnel est simulé.} Les réponses sont
scriptées et l'aiguillage se fait par mots-clés dans les trois langues ; une
phrase qu'aucun mot-clé n'attrape tombe sur la réponse par défaut. Un badge
\textsc{bêta} le signale à l'utilisateur.
\item \textbf{Les traductions arabe et anglaise n'ont pas été relues par un
locuteur natif.} Pour une ligne d'assistance destinée à des enfants en détresse,
le ton compte autant que l'exactitude.
\item \textbf{Les textes juridiques sont incomplets} : les passages qui exigent
un juriste portent un marqueur explicite, et l'écran affiche un bandeau tant
qu'il en reste.
\item \textbf{La mise en page vise le portrait} ; le paysage reste utilisable
sans être soigné.
\item \textbf{Le compte de facturation doit appartenir au \gls{cmrpi}}, non à un
développeur : celui qui le ferme emporte les signalements en cours.
\end{itemize}

\section{Perspectives}

À court terme, la mise en production suppose trois actions qui ne relèvent pas
du code : la relecture des traductions, la rédaction des textes juridiques et le
transfert du projet Firebase au \gls{cmrpi}. Viennent ensuite les extensions
prévues par le cadrage — l'espace destiné aux parents et aux enseignants — et
l'intégration au back-end \gls{emc} et au portail \texttt{e-vigilance}, pour
qu'un signalement mobile alimente directement le système national. L'assistant
conversationnel, enfin, gagnerait à s'appuyer sur un modèle de langue plutôt que
sur des mots-clés, à condition que le traitement reste local ou contractuellement
encadré : les propos d'un enfant en détresse ne sont pas une charge utile
ordinaire.

% ============================================================================
\chapter*{Conclusion}
\addcontentsline{toc}{chapter}{Conclusion}
\markboth{Conclusion}{}
% ============================================================================

Le prototype demandé était modeste ; la question qu'il posait ne l'était pas.
Retirer les comptes utilisateurs, geste évident du point de vue de l'enfant,
retire du même coup le socle sur lequel repose la sécurité usuelle d'une base
\emph{cloud}. Tout le reste — les règles fermées, les fonctions appelables,
l'attestation d'intégrité, l'empreinte du numéro de référence, la conservation
bornée — découle de cette contrainte initiale et non d'un goût pour la
complexité.

Sur le plan personnel, ce stage a été l'occasion de mener un projet de bout en
bout : cadrage, maquettage, développement mobile, conception d'une \gls{api} et
mise en place d'une chaîne d'intégration continue. La leçon la plus durable tient
sans doute à une habitude prise en chemin — écrire, à côté de chaque choix
technique, la raison qui l'a imposé. Un an plus tard, c'est cette trace qui
distingue une décision d'un accident.

Le prototype est fonctionnel et testé. Ce qui lui manque pour servir réellement
ne relève plus de la technique : une relecture humaine des traductions, un
juriste, et un compte au nom du \gls{cmrpi}.

% ============================================================================
%  BIBLIOGRAPHIE, GLOSSAIRE, ANNEXES
% ============================================================================
\cleardoublepage
\printbibliography[heading=bibintoc,title={Bibliographie et références}]

\appendix
\renewcommand{\chaptertitlename}{Annexe}

\chapter{Annexes}
\section{Environnement de développement}
\label{ann:env}

\begin{table}[htbp]
\centering
\small
\caption{Outils et versions}
\label{tab:outils}
\begin{tabularx}{\textwidth}{@{}L{4.2cm} X@{}}
\toprule
\textbf{Outil} & \textbf{Rôle et version} \\
\midrule
Flutter 3.44.6 / Dart 3.12 & Application mobile, version épinglée \\
Node 22, TypeScript & Fonctions \emph{cloud} de 2\ieme{} génération \\
Firebase Firestore & Base de données, région \texttt{europe-west1} \\
Cloud Storage & Preuves, écriture par URL signée uniquement \\
Firebase App Check & Play Integrity / DeviceCheck \\
Émulateurs Firebase & Parcours complet en local, sans facturation \\
Figma & Maquettage des treize écrans \\
Git, GitHub Actions & Versionnement et intégration continue \\
\bottomrule
\end{tabularx}
\end{table}

\begin{lstlisting}[language=bash,caption={Exécution du parcours complet en local},
                   label={lst:local}]
npm run backend                              # emulateurs Firestore, Functions, Storage
flutter run --dart-define=USE_EMULATORS=true # application pointee sur les emulateurs
npm run console                              # console equipe + compte agent local

flutter analyze && flutter test              # 187 tests applicatifs
npm run check                                # types + 103 tests back-end
\end{lstlisting}

\section{Écrans complémentaires de l'application}
\label{ann:ecrans}

\begin{figure}[htbp]
\centering
\begin{subfigure}[b]{0.213\textwidth}\centering
  \phoneshot{\linewidth}{figures/screens/app-09-ressources.png}{App — ressources}
  \caption{Ressources}\end{subfigure}\hfill
\begin{subfigure}[b]{0.213\textwidth}\centering
  \phoneshot{\linewidth}{figures/screens/app-10-contact.png}{App — contact}
  \caption{Contact}\end{subfigure}\hfill
\begin{subfigure}[b]{0.213\textwidth}\centering
  \phoneshot{\linewidth}{figures/screens/app-11-assistant.png}{App — assistant (bêta)}
  \caption{Assistant}\end{subfigure}\hfill
\begin{subfigure}[b]{0.213\textwidth}\centering
  \phoneshot{\linewidth}{figures/screens/app-12-parametres.png}{App — paramètres}
  \caption{Paramètres}\end{subfigure}
\caption[Écrans complémentaires de l'application]{Les trois onglets restants et
l'écran de paramètres, d'où se règlent la langue, les notifications et l'accès
aux textes légaux.}
\label{fig:annexe-ecrans}
\end{figure}

\end{document}
